\chapter{Schr\"odinger's Equation}

%===========================================

%===========================================

\section{Time-Independent Schrödinger Equation}

Separating variables with $\Psi(x,t) = \psi(x)\varphi(t)$:
%
\begin{itemize}
\item Time part:
\[
\frac{d\varphi}{dt} = -\frac{iE}{\hbar}\varphi 
\quad \Rightarrow \quad 
\varphi = \exp\left(-\frac{iEt}{\hbar}\right)
\]
%
\item Spatial part:
\[\boxed{
\hat{H}\psi = E\psi
}\]
\end{itemize}

\section{Infinite Square Well}

For the potential
\[
V(x) = \begin{cases} 0, & 0 \leq x \leq a \\ +\infty, & \text{otherwise} \end{cases}
\]
the Schrödinger equation becomes:
\[
\frac{d^2\psi}{dx^2} = -k^2\psi, \quad
k = \frac{1}{\hbar}\sqrt{2mE}
\]
%
The general solution $\psi = A\sin kx + B\cos kx$ must satisfy boundary conditions $\psi(0) = \psi(a) = 0$.\\
Substituting and normalizing yields the \emph{eigenfunctions and energy levels}:
\[
\psi_n(x) = \sqrt{\frac{2}{a}}\sin\frac{n\pi x}{a}, \quad E_n = \frac{n^2h^2}{8ma^2}
\]
Note that $\psi_n$ forms a complete orthonormal basis for the entire space. Thus any initial state $\Psi(x,0)$ can be expanded via Fourier transform to obtain the time-dependent wave function:
\[
\Psi(x,t) = \sum_n c_n\psi_n\varphi_n, \quad \text{where } c_n = \sqrt{\frac{2}{a}}\int_0^a \sin\frac{n\pi x}{a}\Psi(x,0)\,dx
\]
Since generally $V(x) \to 0$ as $x \to \pm\infty$, we defined two kinds of states to be discussed later on:
\begin{itemize}
    \item $E < 0$: \emph{Bound states} (localized)
    \item $E > 0$: \emph{Scattering states} (transmission/reflection occurs at potential barriers)
\end{itemize}

\section{Delta Function Potential}

For $V(x) = -\alpha\delta(x)$ with $\alpha > 0$:

\subsection{Bound States}
Consider bound states with $E < 0$, the Schrödinger equation is 
\[
\frac{d^2\psi}{dx^2} = \kappa^2\psi, \quad
\kappa = \frac{1}{\hbar}\sqrt{-2mE}
\]
General solution writes $\psi = Ae^{-\kappa x} + Be^{\kappa x}$. \\
For $x < 0$, the $A$ term diverges as $x \to -\infty$, so $A = 0$. Similarly for $x > 0$, the $B$ term must be discarded. 
By continuity of wave function, the remaining coefficients are equal:
\[
\psi(x) = \begin{cases} Be^{\kappa x}, & x < 0 \\ Be^{-\kappa x}, & x > 0 \end{cases}
\]
While the wave function is differentiable everywhere for finite potentials, it is not differentiable at $x = 0$ for the delta potential. Integrating the Schrödinger equation across an infinitesimal interval and taking the limit:
\[
-\frac{\hbar^2}{2m}\int_{-\varepsilon}^{\varepsilon}\frac{d^2\psi}{dx^2}\,dx + \int_{-\varepsilon}^{\varepsilon}V\psi\,dx = E\int_{-\varepsilon}^{\varepsilon}\psi\,dx
\]
As $\varepsilon \to 0$, the right side vanishes (since $\psi$ is finite). Thus:
\[
\lim_{\varepsilon\to 0}\left(\frac{d\psi}{dx}\bigg|_{+\varepsilon} - \frac{d\psi}{dx}\bigg|_{-\varepsilon}\right) = \frac{2m}{\hbar^2}\lim_{\varepsilon\to 0}\int_{-\varepsilon}^{\varepsilon}V\psi\,dx
\]
For finite $V$, the right side vanishes, giving equal derivatives. But for $V(x) = -\alpha\delta(x)$:
\[
\Delta\left(\frac{d\psi}{dx}\right) = -\frac{2m\alpha}{\hbar^2}\psi(0) = -\frac{2m\alpha B}{\hbar^2}
\]
For our solution: 
\[
\left.\frac{d\psi}{dx}\right|_{0^+} \!= -B\kappa, \quad
\left.\frac{d\psi}{dx}\right|_{0^-} \!= B\kappa
\quad\Rightarrow\quad 
\Delta(d\psi/dx) = -2B\kappa
\]
After solving and normalization:
\[
\kappa = \frac{m\alpha}{\hbar^2}, \quad
B = \sqrt{\frac{m\alpha}{\hbar^2}} = \sqrt{\kappa}
\]
Thus the \emph{bound state wave function and energy} are:
\[
\psi(x) = \frac{\sqrt{m\alpha}}{\hbar}\exp\left(-\frac{m\alpha|x|}{\hbar^2}\right), \quad E = -\frac{m\alpha^2}{2\hbar^2}
\]

\subsection{Scattering States}

The Schrödinger equation is then
\[
\frac{d^2\psi}{dx^2} = -k^2\psi, \quad
k = \frac{1}{\hbar}\sqrt{2mE}
\]
The equation solves:
\[\begin{aligned}
\text{For } x < 0&:\ \psi = Ae^{ikx} + Be^{-ikx} \\
\text{For } x > 0&:\ \psi = Fe^{ikx} + Ge^{-ikx}
\end{aligned}\]
Boundary conditions at $x = 0$ are obtained by means similar to bound states:
\begin{align*}
&\text{Continuity: } F + G = A + B \\
&\text{Derivative discontinuity: } ik(F - G - A + B) = -\frac{2m\alpha}{\hbar^2}(A + B)
\end{align*}
Letting $\beta = m\alpha/\hbar^2 k$, we obtain:
\[
F - G = (1 + 2i\beta)A - (1 - 2i\beta)B
\]
Identifying $A$ as incident wave, $B$ as reflected wave, $F$ as transmitted wave, and setting $G = 0$ (no wave from $+\infty$):
\[
B = \frac{i\beta}{1 - i\beta}A, \quad F = \frac{1}{1 - i\beta}A
\]
The \emph{reflection and transmission coefficients} are:
\[
R = \frac{|B|^2}{|A|^2} = \frac{\beta^2}{1 + \beta^2}, \quad T = \frac{|F|^2}{|A|^2} = \frac{1}{1 + \beta^2}
\]
Note that for $V = +\alpha\delta(x)$ with $\alpha > 0$, the bound state disappears, but scattering states remain with finite transmission probability even when the particle energy is insufficient to classically overcome the barrier. This phenomenon of transmission through a classically forbidden region is called \emph{quantum tunneling}.

\section{Finite Square Well}

Consider the potential
\[
V(x) = \begin{cases} -V_0, & -a \leq x \leq a \\ 0, & \text{otherwise} \end{cases} \quad (V_0 > 0)
\]

\subsection{Bound States}

\begin{itemize}

\item For $-a \leq x \leq a$: Schrödinger equation is 
\[
\frac{d^2\psi}{dx^2} = -l^2\psi, \quad
l = \frac{1}{\hbar}\sqrt{2m(E + V_0)}
\]

Solution: $\psi_a(x) = C\sin(lx) + D\cos(lx)$

\item For $|x| > a$: 
\[
\frac{d^2\psi}{dx^2} = \kappa^2\psi, \quad 
\kappa = \frac{1}{\hbar}\sqrt{-2mE}
\]
Solution for $x > a$: $\psi_b(x) = Fe^{-\kappa x}$ (the $e^{\kappa x}$ term is discarded as it diverges at $+\infty$).\\ \\
\item For \emph{even solutions} $\psi_a(x) = D\cos(lx)$:
\begin{align*}
&\text{Continuity: } Fe^{-\kappa a} = D\cos(la) \\
&\text{Differentiability: } -\kappa Fe^{-\kappa a} = -lD\sin(la)
\end{align*}
Dividing $\kappa = l\tan(la)$, then let $z = la$ and $z_0 = \frac{a}{\hbar}\sqrt{2mV_0}$:
\[
\tan z = \sqrt{\left(\frac{z_0}{z}\right)^2 - 1}
\]
\item For \emph{odd solutions}, similar analysis yields:
\[
-\cot z = \sqrt{\left(\frac{z_0}{z}\right)^2 - 1}
\]
\end{itemize}

Therefore the number of bound states is:
\[N=\left\lfloor\frac{2z_0}{\pi}\right\rfloor + 1\]

\subsection{Scattering States}

The wave function has the form:
\[
\psi = \begin{cases} Ae^{ikx} + Be^{-ikx}, & x < -a \\ C\sin(lx) + D\cos(lx), & -a \leq x \leq a \\ Fe^{ikx}, & x > a \end{cases}
\]
The transmission coefficient is:
\[
T^{-1} = 1 + \frac{V_0^2}{4E(E + V_0)}\sin^2\left(\frac{2a}{\hbar}\sqrt{2m(E + V_0)}\right)
\]

\section{Harmonic Oscillator and Ladder Operators}

When a particle is at equilibrium near a potential minimum, Taylor expansion gives:
\[
V(x) = V(x_0) + V'(x_0)(x - x_0) + \frac{1}{2}V''(x_0)(x - x_0)^2 + \cdots
\]
$V(x_0)$ is an irrelevant constant which does not affect dynamics. At equilibrium $V'(x_0) = 0$, so the linear term vanishes. For small oscillations, higher order terms are negligible, leaving only the quadratic term:
\[V(x) \approx \frac{1}{2}V''(x_0)(x - x_0)^2\]
This explains why general potentials near equilibrium can be approximated as harmonic oscillators.\\ \\
For $V(x) = \frac{1}{2}m\omega^2x^2$, the Hamiltonian is:
\[
\hat{H} = \frac{\hat{p}^2}{2m} + \frac{1}{2}m\omega^2\hat{x}^2 = \frac{1}{2m}\left[\hat{p}^2 + (m\omega x)^2\right]
\]
Motivated by the complex factorization $x^2 + y^2 = (ix + y)(-ix + y)$, define the \emph{ladder operators}:
\[
\hat{a}_{\pm} = \frac{1}{\sqrt{2\hbar m\omega}}\left(\mp i\hat{p} + m\omega x\right)
\]
Since $x$ and $p$ do not commute $[\hat{x}, \hat{p}] = i\hbar$, we have:
\[
\hat{a}_-\hat{a}_+ = \frac{1}{2\hbar m\omega}\left[\hat{p}^2 + (m\omega x)^2 - im\omega[\hat{x}, \hat{p}]\right] = \frac{1}{\hbar\omega}\hat{H} + \frac{1}{2}
\]
The operators do not commute: $[\hat{a}_-, \hat{a}_+] = 1$, and they are Hermitian conjugates $\hat{a}_-^\dagger = \hat{a}_+$, thus:
\[\hat{H} = \hbar\omega\left(\hat{a}_-\hat{a}_+ - \frac{1}{2}\right) = \hbar\omega\left(\hat{a}_+\hat{a}_- + \frac{1}{2}\right)\]
If $\psi$ has energy eigenvalue $E$, then $\hat{a}_\pm\psi$ has eigenvalue $E \pm \hbar\omega$:
\begin{align*}
\hat{H}(\hat{a}_+\psi) &= \hbar\omega\left(\hat{a}_+\hat{a}_- + \frac{1}{2}\right)(\hat{a}_+\psi) \\
&= \hbar\omega\left(\hat{a}_+\hat{a}_-\hat{a}_+ + \frac{1}{2}\hat{a}_+\right)\psi \\
&= \hbar\omega\hat{a}_+\left(\hat{a}_-\hat{a}_+ + \frac{1}{2}\right)\psi \\
&= \hat{a}_+(\hat{H} + \hbar\omega)\psi = \hat{a}_+(E + \hbar\omega)\psi = (E + \hbar\omega)(\hat{a}_+\psi)
\end{align*}
Therefore $\hat{a}_\pm$ are called \emph{ladder operators} ($\hat{a}_+$ raises, $\hat{a}_-$ lowers).\\ \\
If we repeatedly apply the lowering operator, the energy would eventually become negative. However, since $\hat{a}_-\psi_0 = 0$ for the ground state, there exists a \emph{minimum energy state} $\psi_0$ satisfying:
\[
\frac{1}{\sqrt{2\hbar m\omega}}\left(\hbar\frac{d}{dx} + m\omega x\right)\psi_0 = 0 \quad \Rightarrow \quad \psi_0 = A\exp\left(-\frac{m\omega x^2}{2\hbar}\right)
\]
Normalizing yields the \emph{ground state wave function and zero-point energy}:
\[
\psi_0(x) = \left(\frac{m\omega}{\pi\hbar}\right)^{1/4}\exp\left(-\frac{m\omega}{2\hbar}x^2\right), \quad E_0 = \frac{1}{2}\hbar\omega
\]
All excited states are obtained by applying the raising operator: $\psi_n = A_n(\hat{a}_+)^n\psi_0$, which give rise to the quantization of energies.\\ \\
To calculate $A_n$, note that $\hat{a}_+\psi_n$ is proportional to $\psi_{n+1}$ and $\hat{a}_-\psi_{n+1}$ is proportional to $\psi_n$. Since $\hat{a}_-\hat{a}_+\psi_n = (n+1)\psi_n$ which can be derived from $[\hat{a}_-, \hat{a}_+] = 1$ and $\hat{H} = \hbar\omega(\hat{a}_-\hat{a}_+ - \frac{1}{2})$, using the Hermitian conjugate property gives: (see \ref{calc})
\[
\langle\hat{a}_+\psi_n|\hat{a}_+\psi_n\rangle = \langle\psi_n|\hat{a}_-\hat{a}_+\psi_n\rangle = (n+1)\langle\psi_n|\psi_n\rangle = n+1
\]
Thus $|c_n|^2 = n+1$ where $\hat{a}_+\psi_n = c_n\psi_{n+1}$. Similarly for $\hat{a}_-$, we obtain:
\[
\hat{a}_+\psi_n = \sqrt{n+1}\psi_{n+1}, \quad \hat{a}_-\psi_n = \sqrt{n}\psi_{n-1}
\]
Therefore we obtain the wave function and energy of harmonic oscillators:
\[
\psi_n = \frac{1}{\sqrt{n!}}(\hat{a}_+)^n\psi_0, \quad E_n = \left(n + \frac{1}{2}\right)\hbar\omega
\]